% File: ai-nhap/2026-11-15_16_de_tham_khao_phu_nhuan_de_1_va_phu_nhuan_de_2.tex
% Nguồn: Tuyển tập đề tham khảo môn Toán Tuyển sinh 10 TP.HCM 2025-2026
% Gồm 14 bài: Quận Phú Nhuận - Đề 1 (7 bài) và Quận Phú Nhuận - Đề 2 (7 bài)

% ==============================================================================
% PHẦN 1: ĐỀ THAM KHẢO TUYỂN SINH 10 QUẬN PHÚ NHUẬN – ĐỀ 1 (7 BÀI)
% ==============================================================================

\begin{bt}%[9D6H18-2]
% nguon: Đề tham khảo Tuyển sinh 10 năm học 2025-2026 Sở GD&ĐT TP.HCM – Quận Phú Nhuận - Đề 1 (Bài 1)
Cho hàm số bậc hai $y = 2x^2$ có đồ thị là đường parabol $(P)$.
\begin{enumerate}[a)]
    \item Vẽ đồ thị $(P)$ trên mặt phẳng tọa độ $Oxy$.
    \item Tìm tọa độ các điểm thuộc parabol $(P)$ có tung độ bằng $2$.
\end{enumerate}
\loigiai{
\begin{enumerate}[a)]
    \item Bảng giá trị tương ứng của hàm số $y = 2x^2$:
    \begin{center}
    \begin{tabular}{|c|c|c|c|c|c|}
        \hline
        $x$ & $-2$ & $-1$ & $0$ & $1$ & $2$ \\
        \hline
        $y = 2x^2$ & $8$ & $2$ & $0$ & $2$ & $8$ \\
        \hline
    \end{tabular}
    \end{center}
    Đồ thị $(P)$ là một parabol có đỉnh tại gốc tọa độ $O(0; 0)$, đối xứng qua trục tung $Oy$ và bề lõm quay lên trên.
    \item Các điểm thuộc $(P)$ có tung độ bằng $2$ thỏa mãn phương trình:
    $y = 2 \iff 2x^2 = 2 \iff x^2 = 1 \iff x = \pm 1$.
    Với $x = 1 \implies y = 2$, ta có điểm $M_1(1; 2)$.
    Với $x = -1 \implies y = 2$, ta có điểm $M_2(-1; 2)$.
    Vậy có hai điểm thỏa mãn yêu cầu là $(1; 2)$ và $(-1; 2)$.
\end{enumerate}
}
\end{bt}

\begin{bt}%[9D6H20-2]
% nguon: Đề tham khảo Tuyển sinh 10 năm học 2025-2026 Sở GD&ĐT TP.HCM – Quận Phú Nhuận - Đề 1 (Bài 2)
Cho phương trình bậc hai một ẩn $4x^2 - 3x - 2 = 0$.
\begin{enumerate}[a)]
    \item Chứng minh phương trình luôn có hai nghiệm phân biệt $x_1, x_2$.
    \item Không giải phương trình, hãy tính giá trị của biểu thức $A = (2x_1 - 3)(2x_2 - 3)$.
\end{enumerate}
\loigiai{
\begin{enumerate}[a)]
    \item Phương trình $4x^2 - 3x - 2 = 0$ có hệ số $a = 4$, $c = -2$.
    Vì $a \cdot c = 4 \cdot (-2) = -8 < 0$ nên phương trình luôn có hai nghiệm phân biệt trái dấu $x_1, x_2$.
    \item Áp dụng định lí Viète cho phương trình $4x^2 - 3x - 2 = 0$, ta có:
    $\begin{cases} S = x_1 + x_2 = -\dfrac{b}{a} = \dfrac{3}{4} \\ P = x_1 x_2 = \dfrac{c}{a} = -\dfrac{2}{4} = -\dfrac{1}{2}. \end{cases}$
    Khai triển và biến đổi biểu thức $A$:
    $A = (2x_1 - 3)(2x_2 - 3) = 4x_1 x_2 - 6(x_1 + x_2) + 9 = 4P - 6S + 9$.
    Thay các giá trị $S = \dfrac{3}{4}$ và $P = -\dfrac{1}{2}$ vào, ta được:
    $A = 4 \cdot \left(-\dfrac{1}{2}\right) - 6 \cdot \dfrac{3}{4} + 9 = -2 - \dfrac{9}{2} + 9 = 7 - \dfrac{9}{2} = \dfrac{5}{2} = 2{,}5$.
\end{enumerate}
}
\end{bt}

\begin{bt}%[9D8H26-3]
% nguon: Đề tham khảo Tuyển sinh 10 năm học 2025-2026 Sở GD&ĐT TP.HCM – Quận Phú Nhuận - Đề 1 (Bài 3)
Bạn Minh gieo một đồng xu cân đối, đồng chất và bạn Phương rút ngẫu nhiên một tấm thẻ từ một chiếc hộp chứa $5$ tấm thẻ được đánh số từ $1$ đến $5$. Tính xác suất của mỗi biến cố sau:
\begin{enumerate}[a)]
    \item Biến cố $A$: ``Rút được tấm thẻ ghi số lẻ''.
    \item Biến cố $B$: ``Rút được tấm thẻ ghi số chẵn và đồng xu xuất hiện mặt sấp''.
    \item Biến cố $C$: ``Rút được tấm thẻ ghi số $5$ hoặc đồng xu xuất hiện mặt ngửa''.
\end{enumerate}
\loigiai{
Đồng xu có $2$ khả năng (mặt Sấp $S$ hoặc mặt Ngửa $N$), hộp thẻ có $5$ khả năng $\{1; 2; 3; 4; 5\}$.
Không gian mẫu gồm tập hợp các cặp $(k; m)$ với $k \in \{1; 2; 3; 4; 5\}$ và $m \in \{S; N\}$:
Tổng số kết quả có thể xảy ra là $n(\Omega) = 5 \cdot 2 = 10$.
\begin{enumerate}[a)]
    \item Các tấm thẻ ghi số lẻ là $\{1; 3; 5\}$ ($3$ tấm), kết hợp với $2$ mặt của đồng xu:
    Số kết quả thuận lợi cho biến cố $A$ là: $n(A) = 3 \cdot 2 = 6$.
    Xác suất của biến cố $A$ là:
    $P(A) = \dfrac{n(A)}{n(\Omega)} = \dfrac{6}{10} = \dfrac{3}{5} = 0{,}6$.
    \item Các tấm thẻ ghi số chẵn là $\{2; 4\}$ ($2$ tấm) và đồng xu mặt sấp ($S$):
    Các kết quả thuận lợi cho biến cố $B$ là: $\{(2; S); (4; S)\} \implies n(B) = 2$.
    Xác suất của biến cố $B$ là:
    $P(B) = \dfrac{n(B)}{n(\Omega)} = \dfrac{2}{10} = \dfrac{1}{5} = 0{,}2$.
    \item Biến cố $C$ xảy ra khi rút được thẻ số $5$ hoặc đồng xu mặt ngửa:
    Các kết quả có thẻ số $5$: $(5; S)$ và $(5; N)$ ($2$ kết quả).
    Các kết quả có đồng xu mặt ngửa: $(1; N), (2; N), (3; N), (4; N), (5; N)$ ($5$ kết quả).
    Hợp của hai tập kết quả gồm $6$ phần tử riêng biệt:
    $\{(1; N); (2; N); (3; N); (4; N); (5; N); (5; S)\} \implies n(C) = 6$.
    Xác suất của biến cố $C$ là:
    $P(C) = \dfrac{n(C)}{n(\Omega)} = \dfrac{6}{10} = \dfrac{3}{5} = 0{,}6$.
\end{enumerate}
}
\end{bt}

\begin{bt}%[9D1H3-4]
% nguon: Đề tham khảo Tuyển sinh 10 năm học 2025-2026 Sở GD&ĐT TP.HCM – Quận Phú Nhuận - Đề 1 (Bài 4)
Để kích cầu tiêu dùng dịp nghỉ lễ 30/4, một siêu thị điện máy thực hiện đợt giảm giá lớn: tivi được giảm $40\%$ so với giá niêm yết và máy giặt được giảm $25\%$ so với giá niêm yết. Tổng giá niêm yết của một chiếc tivi và một chiếc máy giặt là $25{,}4$ triệu đồng. Bác Hai đã mua một chiếc tivi và một chiếc máy giặt trong dịp này với tổng số tiền thực tế phải trả là $16{,}77$ triệu đồng. Hỏi giá niêm yết ban đầu của mỗi món hàng khi chưa giảm giá là bao nhiêu?
\loigiai{
Gọi giá niêm yết của chiếc tivi là $x$ và của chiếc máy giặt là $y$ ($x, y > 0$, đơn vị: triệu đồng).
Vì tổng giá niêm yết của hai món hàng là $25{,}4$ triệu đồng nên ta có phương trình:
$x + y = 25{,}4$. \quad (1)
Giá bán thực tế của chiếc tivi sau khi giảm $40\%$ là:
$x \cdot (1 - 0{,}40) = 0{,}6x$ (triệu đồng).
Giá bán thực tế của chiếc máy giặt sau khi giảm $25\%$ là:
$y \cdot (1 - 0{,}25) = 0{,}75y$ (triệu đồng).
Tổng số tiền thực tế bác Hai đã thanh toán là $16{,}77$ triệu đồng nên ta có phương trình:
$0{,}6x + 0{,}75y = 16{,}77$. \quad (2)
Từ (1) và (2) ta có hệ phương trình:
$\begin{cases} x + y = 25{,}4 \\ 0{,}6x + 0{,}75y = 16{,}77 \end{cases} \iff \begin{cases} 0{,}6x + 0{,}6y = 15{,}24 \\ 0{,}6x + 0{,}75y = 16{,}77 \end{cases} \iff \begin{cases} 0{,}15y = 1{,}53 \\ x = 25{,}4 - y \end{cases} \iff \begin{cases} y = 10{,}2 \\ x = 15{,}2 \end{cases}$ (thỏa mãn điều kiện).
Vậy giá niêm yết ban đầu của chiếc tivi là $15{,}2$ triệu đồng ($15\,200\,000$ đồng) và của chiếc máy giặt là $10{,}2$ triệu đồng ($10\,200\,000$ đồng).
}
\end{bt}

\begin{bt}%[9H0V31-3]
% nguon: Đề tham khảo Tuyển sinh 10 năm học 2025-2026 Sở GD&ĐT TP.HCM – Quận Phú Nhuận - Đề 1 (Bài 5)
Để làm một chiếc quạt giấy thủ công, bạn Phương dùng các nan tre dài $30\text{ cm}$ chốt lại ở một đầu làm tâm quay $O$. Bạn cắt giấy màu thành hình vành quạt giới hạn bởi hai hình quạt tròn đồng tâm: hình quạt lớn $OAB$ bán kính $R = 30\text{ cm}$ và hình quạt nhỏ $OCD$ bán kính $r = 10\text{ cm}$, cùng có góc ở tâm là $120^\circ$. Phần giấy vành quạt này được dán phủ kín lên cả hai mặt của các nan quạt. Hãy tính tổng diện tích giấy thủ công cần dùng (kết quả làm tròn đến hàng phần trăm của $\text{cm}^2$). Lấy $\pi \approx 3{,}1416$.
\loigiai{
Diện tích của một mặt giấy hình vành quạt bằng hiệu diện tích giữa hình quạt tròn lớn và hình quạt tròn nhỏ:
$S_1 = \dfrac{\pi R^2 \cdot 120^\circ}{360^\circ} - \dfrac{\pi r^2 \cdot 120^\circ}{360^\circ} = \dfrac{\pi}{3}(R^2 - r^2) = \dfrac{\pi}{3}(30^2 - 10^2) = \dfrac{800\pi}{3}\text{ (cm}^2\text{)}$.
Vì quạt được dán giấy ở cả hai mặt (mặt trước và mặt sau) nên tổng diện tích giấy thủ công cần dùng là:
$S = 2 \cdot S_1 = 2 \cdot \dfrac{800\pi}{3} = \dfrac{1600\pi}{3} \approx 1675{,}52\text{ cm}^2$.
Vậy tổng diện tích giấy cần dùng xấp xỉ $1675{,}52\text{ cm}^2$.
}
\end{bt}

\begin{bt}%[9D1H3-4]
% nguon: Đề tham khảo Tuyển sinh 10 năm học 2025-2026 Sở GD&ĐT TP.HCM – Quận Phú Nhuận - Đề 1 (Bài 6)
Một khách sạn tại thành phố Đà Lạt áp dụng mức giá thuê phòng trong dịp lễ Tết như sau: đơn giá thuê phòng là $300\,000$ đồng mỗi buổi và chịu thuế giá trị gia tăng GTGT là $8\%$. Ngoài ra, khách sạn thu thêm một khoản phụ thu dịch vụ lễ Tết cố định là $50\,000$ đồng cho mỗi phòng (chỉ thu một lần duy nhất cho toàn bộ thời gian lưu trú).
\begin{enumerate}[a)]
    \item Gọi $x$ là số buổi khách lưu trú và $y$ là tổng số tiền khách phải thanh toán (đồng). Hãy lập biểu thức biểu diễn $y$ theo $x$.
    \item Anh An đến thuê phòng tại khách sạn này và khi thanh toán đã trả tổng số tiền là $1\,346\,000$ đồng. Hỏi anh An đã lưu trú tại khách sạn trong bao nhiêu buổi?
\end{enumerate}
\loigiai{
\begin{enumerate}[a)]
    \item Đơn giá mỗi buổi thuê phòng sau khi đã tính thuế GTGT $8\%$ là:
    $300\,000 \cdot (1 + 0{,}08) = 324\,000$ (đồng/buổi).
    Biểu thức biểu diễn tổng số tiền $y$ phải trả theo số buổi lưu trú $x$ là:
    $y = 324\,000x + 50\,000$ (đồng).
    \item Với tổng số tiền thanh toán là $1\,346\,000$ đồng, ta có phương trình:
    $324\,000x + 50\,000 = 1\,346\,000 \iff 324\,000x = 1\,296\,000 \iff x = 4$ (buổi).
    Vậy anh An đã lưu trú tại khách sạn trong $4$ buổi.
\end{enumerate}
}
\end{bt}

\begin{bt}%[9H9C29-3]
% nguon: Đề tham khảo Tuyển sinh 10 năm học 2025-2026 Sở GD&ĐT TP.HCM – Quận Phú Nhuận - Đề 1 (Bài 7)
Cho đường tròn $(O; R)$ có đường kính $AC$ và đường thẳng $d$ là tiếp tuyến của $(O)$ tại điểm $A$. Trên đường thẳng $d$ lấy điểm $M$ sao cho $AM > AO$. Từ $M$ vẽ tiếp tuyến thứ hai $MB$ với $(O)$ ($B$ là tiếp điểm, $B$ khác $A$).
\begin{enumerate}[a)]
    \item Chứng minh bốn điểm $M, A, O, B$ cùng thuộc một đường tròn và $OM \perp AB$.
    \item Gọi $D$ là giao điểm của đoạn thẳng $MO$ với đường tròn $(O)$. Tia $AD$ cắt đoạn thẳng $MB$ tại $E$. Chứng minh $EB^2 = EA \cdot ED$.
    \item Đường phân giác của góc $\widehat{ABC}$ trong tam giác $ABC$ cắt đường tròn $(O; R)$ tại điểm $K$ ($K$ khác $B$). Kẻ $BI \perp AC$ tại $I$. Đặt $BI = x$ (với $x < R$). Hãy tính diện tích tam giác $BIK$ theo $R$ và $x$.
\end{enumerate}
\loigiai{
\begin{enumerate}[a)]
    \item Vì $MA, MB$ là hai tiếp tuyến của $(O)$ nên $\widehat{MAO} = 90^\circ$ và $\widehat{MBO} = 90^\circ$.
    Tứ giác $MAOB$ có hai góc đối vuông $\widehat{MAO} + \widehat{MBO} = 180^\circ$ nên bốn điểm $M, A, O, B$ cùng thuộc đường tròn đường kính $MO$.
    Lại có $MA = MB$ và $OA = OB = R$, suy ra $OM$ là đường trung trực của đoạn thẳng $AB$, do đó $OM \perp AB$.
    \item Vì $D \in (O)$ và $OM$ là trục đối xứng của dây cung $AB$ nên cung $DA$ bằng cung $DB$, suy ra $\widehat{EBD} = \widehat{EAB}$ (góc tạo bởi tiếp tuyến và dây cung bằng góc nội tiếp cùng chắn hai cung bằng nhau).
    Xét $\triangle EBD$ và $\triangle EAB$, ta có góc $\widehat{AEB}$ chung và $\widehat{EBD} = \widehat{EAB}$, suy ra $\triangle EBD \backsim \triangle EAB$ (g.g).
    Do đó $\dfrac{EB}{EA} = \dfrac{ED}{EB} \implies EB^2 = EA \cdot ED$ (đpcm).
    \item Trong tam giác vuông $ABC$ nội tiếp chắn nửa đường tròn đường kính $AC$, ta có $\widehat{ABC} = 90^\circ$.
    Vì $BK$ là tia phân giác của góc vuông $\widehat{ABC}$ nên $\widehat{ABK} = \widehat{CBK} = 45^\circ$.
    Suy ra $K$ là điểm chính giữa của cung $AC$, do đó $OK \perp AC$ và $OK = R$.
    Khoảng cách từ $K$ đến $AC$ bằng $R$ và từ $B$ đến $AC$ bằng $BI = x$.
    Diện tích tam giác $BIK$ được tính theo hệ thức hình học là:
    $S_{\triangle BIK} = \dfrac{1}{2} x \cdot R$.
\end{enumerate}
}
\end{bt}


% ==============================================================================
% PHẦN 2: ĐỀ THAM KHẢO TUYỂN SINH 10 QUẬN PHÚ NHUẬN – ĐỀ 2 (7 BÀI)
% ==============================================================================

\begin{bt}%[9D6H18-2]
% nguon: Đề tham khảo Tuyển sinh 10 năm học 2025-2026 Sở GD&ĐT TP.HCM – Quận Phú Nhuận - Đề 2 (Bài 1)
Cho hàm số bậc hai $y = \dfrac{1}{2}x^2$ có đồ thị là đường parabol $(P)$.
\begin{enumerate}[a)]
    \item Vẽ đồ thị $(P)$ trên mặt phẳng tọa độ $Oxy$.
    \item Bạn Minh đang thiết kế một trò chơi điện tử mô phỏng quỹ đạo bay của một quả bóng theo đường parabol $(P)\colon y = \dfrac{1}{2}x^2$. Minh muốn tìm các vị trí trên quỹ đạo mà độ cao của quả bóng so với mặt đất đạt đúng $0{,}5\text{ m}$. Hỏi tại những điểm nào trên mặt phẳng tọa độ thì quả bóng đạt được độ cao này?
\end{enumerate}
\loigiai{
\begin{enumerate}[a)]
    \item Bảng giá trị tương ứng của hàm số $y = \dfrac{1}{2}x^2$:
    \begin{center}
    \begin{tabular}{|c|c|c|c|c|c|}
        \hline
        $x$ & $-4$ & $-2$ & $0$ & $2$ & $4$ \\
        \hline
        $y = \dfrac{1}{2}x^2$ & $8$ & $2$ & $0$ & $2$ & $8$ \\
        \hline
    \end{tabular}
    \end{center}
    Đồ thị $(P)$ là một parabol có đỉnh $O(0; 0)$, đối xứng qua trục tung $Oy$ và bề lõm quay lên trên.
    \item Vị trí quả bóng đạt độ cao $0{,}5\text{ m}$ ứng với tung độ $y = 0{,}5$.
    Thay vào phương trình của $(P)$, ta có:
    $\dfrac{1}{2}x^2 = 0{,}5 \iff x^2 = 1 \iff x = \pm 1$.
    Vậy trên mặt phẳng tọa độ có hai vị trí mà quả bóng đạt được độ cao $0{,}5\text{ m}$ là điểm $A(1; 0{,}5)$ và điểm $B(-1; 0{,}5)$.
\end{enumerate}
}
\end{bt}

\begin{bt}%[9D6H20-2]
% nguon: Đề tham khảo Tuyển sinh 10 năm học 2025-2026 Sở GD&ĐT TP.HCM – Quận Phú Nhuận - Đề 2 (Bài 2)
Cho phương trình bậc hai $2x^2 - x - 2 = 0$ có hai nghiệm phân biệt $x_1, x_2$. Không giải phương trình, hãy tính giá trị của biểu thức $P = \sqrt{3 - x_1^2 - x_2^2}$.
\loigiai{
Phương trình $2x^2 - x - 2 = 0$ có $a \cdot c = 2 \cdot (-2) = -4 < 0$ nên luôn có hai nghiệm phân biệt trái dấu $x_1, x_2$.
Theo định lí Viète, ta có:
$\begin{cases} S = x_1 + x_2 = -\dfrac{b}{a} = \dfrac{1}{2} \\ P_0 = x_1 x_2 = \dfrac{c}{a} = -\dfrac{2}{2} = -1. \end{cases}$
Tổng bình phương hai nghiệm là:
$x_1^2 + x_2^2 = (x_1 + x_2)^2 - 2x_1 x_2 = S^2 - 2P_0 = \left(\dfrac{1}{2}\right)^2 - 2 \cdot (-1) = \dfrac{1}{4} + 2 = \dfrac{9}{4}$.
Thay vào biểu thức $P$, ta được:
$P = \sqrt{3 - (x_1^2 + x_2^2)} = \sqrt{3 - \dfrac{9}{4}} = \sqrt{\dfrac{3}{4}} = \dfrac{\sqrt{3}}{2}$.
}
\end{bt}

\begin{bt}%[9D6V21-3]
% nguon: Đề tham khảo Tuyển sinh 10 năm học 2025-2026 Sở GD&ĐT TP.HCM – Quận Phú Nhuận - Đề 2 (Bài 3)
Khu vườn trồng hoa hình chữ nhật của gia đình bác An có kích thước ban đầu là chiều dài $50\text{ m}$ và chiều rộng $30\text{ m}$. Bác An dự định mở rộng thêm diện tích vườn bằng cách nới rộng thêm mỗi chiều một đoạn bằng $x$ mét về phía ngoài ($x > 0$).
\begin{enumerate}[a)]
    \item Viết biểu thức $S$ biểu diễn diện tích vườn hoa hình chữ nhật sau khi mở rộng theo $x$.
    \item Biết rằng phần diện tích tăng thêm của vườn sau khi mở rộng là $1050\text{ m}^2$. Tìm giá trị của $x$ (kết quả làm tròn đến hàng phần mười của mét).
\end{enumerate}
\loigiai{
\begin{enumerate}[a)]
    \item Chiều dài và chiều rộng của vườn hoa sau khi mở rộng lần lượt là $50 + x$ (m) và $30 + x$ (m).
    Biểu thức diện tích vườn hoa sau khi mở rộng là:
    $S = (50 + x)(30 + x) = x^2 + 80x + 1500\text{ (m}^2\text{)}$.
    \item Diện tích vườn hoa ban đầu là: $50 \cdot 30 = 1500\text{ m}^2$.
    Phần diện tích tăng thêm là:
    $S_{\text{tăng}} = (x^2 + 80x + 1500) - 1500 = x^2 + 80x\text{ (m}^2\text{)}$.
    Theo đề bài, phần diện tích tăng thêm là $1050\text{ m}^2$, ta có phương trình:
    $x^2 + 80x = 1050 \iff x^2 + 80x - 1050 = 0$.
    Ta có $\Delta' = 40^2 - 1 \cdot (-1050) = 1600 + 1050 = 2650 > 0$.
    Vì $x > 0$ nên nghiệm thỏa mãn là:
    $x = -40 + \sqrt{2650} \approx 11{,}5\text{ m}$.
    Vậy giá trị của $x$ xấp xỉ $11{,}5\text{ m}$.
\end{enumerate}
}
\end{bt}

\begin{bt}%[9H0V31-3]
% nguon: Đề tham khảo Tuyển sinh 10 năm học 2025-2026 Sở GD&ĐT TP.HCM – Quận Phú Nhuận - Đề 2 (Bài 4)
Nước giải khát thường đựng trong lon nhôm. Cỡ lon truyền thống chứa khoảng $330\text{ ml}$ dung dịch, có dạng hình trụ với chiều cao $10{,}2\text{ cm}$ và đường kính đáy $6{,}42\text{ cm}$. Hiện nay, nhiều nhà sản xuất thiết kế mẫu lon thon cao mới có chiều cao $13{,}41\text{ cm}$ và đường kính đáy $5{,}6\text{ cm}$. Lấy $\pi \approx 3{,}14$.
\begin{enumerate}[a)]
    \item Lon nước ngọt mẫu thon cao có chứa được trọn vẹn $330\text{ ml}$ dung dịch của lon truyền thống không? Vì sao?
    \item Hãy so sánh diện tích toàn phần của hai mẫu lon và giải thích vì sao chi phí sản xuất vỏ lon mẫu thon cao lại tốn kém hơn mẫu lon truyền thống.
\end{enumerate}
\loigiai{
\begin{enumerate}[a)]
    \item Bán kính đáy của lon mẫu thon cao là:
    $r_2 = \dfrac{5{,}6}{2} = 2{,}8\text{ cm}$.
    Thể tích của lon mẫu thon cao là:
    $V_2 = \pi r_2^2 h_2 \approx 3{,}14 \cdot (2{,}8)^2 \cdot 13{,}41 \approx 330{,}12\text{ cm}^3 = 330{,}12\text{ ml}$.
    Vì $330{,}12\text{ ml} > 330\text{ ml}$ nên lon mẫu thon cao hoàn toàn chứa được hết lượng nước ngọt của lon truyền thống.
    \item Bán kính đáy của lon truyền thống là $r_1 = \dfrac{6{,}42}{2} = 3{,}21\text{ cm}$.
    Diện tích toàn phần của lon truyền thống là:
    $S_{\text{tp1}} = 2\pi r_1 h_1 + 2\pi r_1^2 \approx 2 \cdot 3{,}14 \cdot 3{,}21 \cdot 10{,}2 + 2 \cdot 3{,}14 \cdot (3{,}21)^2 \approx 205{,}62 + 64{,}71 = 270{,}33\text{ cm}^2$.
    Diện tích toàn phần của lon mẫu thon cao là:
    $S_{\text{tp2}} = 2\pi r_2 h_2 + 2\pi r_2^2 \approx 2 \cdot 3{,}14 \cdot 2{,}8 \cdot 13{,}41 + 2 \cdot 3{,}14 \cdot (2{,}8)^2 \approx 235{,}79 + 49{,}24 = 285{,}03\text{ cm}^2$.
    Vì $S_{\text{tp2}} > S_{\text{tp1}}$ nên lượng nhôm dùng để chế tạo vỏ lon mẫu thon cao tốn nhiều hơn so với vỏ lon truyền thống, dẫn đến chi phí sản xuất tốn kém hơn.
\end{enumerate}
}
\end{bt}

\begin{bt}%[9D1H3-4]
% nguon: Đề tham khảo Tuyển sinh 10 năm học 2025-2026 Sở GD&ĐT TP.HCM – Quận Phú Nhuận - Đề 2 (Bài 5)
Trong kỳ thi tuyển sinh lớp 10, trường THPT A tuyển học sinh qua ba đợt nguyện vọng:
\begin{itemize}
    \item Ở nguyện vọng 1, số học sinh nam được chia đều thành các nhóm $18$ bạn, học sinh nữ chia đều thành các nhóm $15$ bạn thì vừa đủ.
    \item Ở nguyện vọng 2, trường tuyển thêm $112$ bạn; khi đó số nhóm nam tăng thêm $2$ nhóm và mỗi nhóm nam có thêm $2$ bạn; số nhóm nữ tăng thêm $1$ nhóm và mỗi nhóm nữ có thêm $3$ bạn.
    \item Ở nguyện vọng 3, trường tuyển thêm $86$ bạn và chia đều vừa đủ vào các nhóm hiện có: mỗi nhóm nam nhận thêm $3$ bạn và mỗi nhóm nữ nhận thêm $4$ bạn.
\end{itemize}
Hãy tính tổng số học sinh đã trúng tuyển vào trường THPT A và tỉ lệ phần trăm học sinh nam, học sinh nữ trúng tuyển.
\loigiai{
Gọi số nhóm học sinh nam và số nhóm học sinh nữ ở nguyện vọng 1 lần lượt là $x$ và $y$ ($x, y \in \mathbb{N}^*$).
Ở nguyện vọng 1: số nam là $18x$, số nữ là $15y$.
Ở nguyện vọng 2:
Số nhóm nam là $x + 2$, mỗi nhóm có $18 + 2 = 20$ bạn. Số nam tăng thêm là:
$(x + 2) \cdot 20 - 18x = 2x + 40$.
Số nhóm nữ là $y + 1$, mỗi nhóm có $15 + 3 = 18$ bạn. Số nữ tăng thêm là:
$(y + 1) \cdot 18 - 15y = 3y + 18$.
Tổng số học sinh tuyển thêm ở nguyện vọng 2 là $112$ bạn nên ta có phương trình:
$(2x + 40) + (3y + 18) = 112 \iff 2x + 3y = 54$. \quad (1)
Ở nguyện vọng 3: trường tuyển thêm $86$ học sinh chia vào $x + 2$ nhóm nam (mỗi nhóm thêm $3$ bạn) và $y + 1$ nhóm nữ (mỗi nhóm thêm $4$ bạn):
$3(x + 2) + 4(y + 1) = 86 \iff 3x + 4y + 10 = 86 \iff 3x + 4y = 76$. \quad (2)
Từ (1) và (2) ta có hệ phương trình:
$\begin{cases} 2x + 3y = 54 \\ 3x + 4y = 76 \end{cases} \iff \begin{cases} 6x + 9y = 162 \\ 6x + 8y = 152 \end{cases} \iff \begin{cases} y = 10 \\ x = 12 \end{cases}$ (thỏa mãn điều kiện).
Số học sinh nam trúng tuyển sau ba đợt là:
$N_{\text{nam}} = 18 \cdot 12 + 2 \cdot 12 + 40 + 3 \cdot (12 + 2) = 216 + 64 + 42 = 322$ (học sinh).
Số học sinh nữ trúng tuyển sau ba đợt là:
$N_{\text{nữ}} = 15 \cdot 10 + 3 \cdot 10 + 18 + 4 \cdot (10 + 1) = 150 + 48 + 44 = 242$ (học sinh).
Tổng số học sinh trúng tuyển vào trường THPT A là:
$322 + 242 = 564$ (học sinh).
Tỉ lệ phần trăm học sinh nam trúng tuyển là: $\dfrac{322}{564} \cdot 100\% \approx 57{,}09\%$.
Tỉ lệ phần trăm học sinh nữ trúng tuyển là: $\dfrac{242}{564} \cdot 100\% \approx 42{,}91\%$.
}
\end{bt}

\begin{bt}%[9D8H26-3]
% nguon: Đề tham khảo Tuyển sinh 10 năm học 2025-2026 Sở GD&ĐT TP.HCM – Quận Phú Nhuận - Đề 2 (Bài 6)
Thống kê kết quả kiểm tra môn Tiếng Anh của học sinh lớp 9A được ghi lại theo tỉ lệ phần trăm như sau: có $20\%$ học sinh đạt $7$ điểm, $40\%$ học sinh đạt $8$ điểm, $30\%$ học sinh đạt $9$ điểm và $10\%$ học sinh đạt $10$ điểm. Biết rằng có đúng $4$ bạn đạt điểm $10$.
\begin{enumerate}[a)]
    \item Hỏi lớp 9A có tất cả bao nhiêu học sinh và có bao nhiêu bạn đạt điểm khá ($7$ điểm)?
    \item Có bao nhiêu học sinh đạt trên $8$ điểm? Tính xác suất của biến cố $A$: ``Học sinh được chọn ngẫu nhiên trong lớp đạt điểm tốt (từ $9$ điểm trở lên)''.
\end{enumerate}
\loigiai{
\begin{enumerate}[a)]
    \item Số học sinh đạt điểm $10$ chiếm $10\%$ tổng số học sinh của lớp.
    Tổng số học sinh của lớp 9A là:
    $4 : 10\% = 40$ (học sinh).
    Số học sinh đạt điểm khá ($7$ điểm) là:
    $40 \cdot 20\% = 8$ (học sinh).
    \item Số học sinh đạt trên $8$ điểm gồm các bạn đạt điểm $9$ và điểm $10$, chiếm tỉ lệ:
    $30\% + 10\% = 40\%$.
    Số học sinh đạt trên $8$ điểm là:
    $40 \cdot 40\% = 16$ (học sinh).
    Xác suất để chọn ngẫu nhiên một bạn đạt điểm tốt là:
    $P(A) = \dfrac{16}{40} = 0{,}4$.
\end{enumerate}
}
\end{bt}

\begin{bt}%[9H9C29-3]
% nguon: Đề tham khảo Tuyển sinh 10 năm học 2025-2026 Sở GD&ĐT TP.HCM – Quận Phú Nhuận - Đề 2 (Bài 7)
Cho điểm $S$ nằm bên ngoài đường tròn $(O; R)$ sao cho đoạn thẳng $SO = 2R$. Qua $S$ vẽ hai tiếp tuyến $SA$ và $SB$ với đường tròn $(O)$ ($A, B$ là các tiếp điểm). Gọi $I$ là giao điểm của đoạn thẳng $SO$ và dây cung $AB$.
\begin{enumerate}[a)]
    \item Chứng minh bốn điểm $S, A, O, B$ cùng thuộc một đường tròn và $SO \perp AB$ tại $I$.
    \item Vẽ đường kính $AD$ của đường tròn $(O)$, đoạn thẳng $SD$ cắt $(O)$ tại điểm thứ hai là $E$ ($E$ khác $D$). Chứng minh $SE \cdot SD = SA^2$ và $SE \cdot SD = SI \cdot SO$.
    \item Cho biết bán kính $R = 8\text{ cm}$. Hãy tính diện tích hình phẳng giới hạn bởi hai bán kính $OA, OB$ và cung nhỏ $AB$ (kết quả làm tròn đến hàng đơn vị của $\text{cm}^2$). Lấy $\pi \approx 3{,}1416$.
\end{enumerate}
\loigiai{
\begin{enumerate}[a)]
    \item Vì $SA, SB$ là hai tiếp tuyến của $(O)$ nên $\widehat{SAO} = 90^\circ$ và $\widehat{SBO} = 90^\circ$.
    Tứ giác $SAOB$ có hai góc đối vuông nên nội tiếp đường tròn đường kính $SO$.
    Lại có $SA = SB$ và $OA = OB = R$, suy ra $SO$ là đường trung trực của $AB$, do đó $SO \perp AB$ tại $I$.
    \item Vì $AD$ là đường kính của $(O)$ và $E \in (O)$ nên $\widehat{AED} = 90^\circ$, suy ra $AE \perp SD$ tại $E$.
    Tam giác vuông $SAD$ tại $A$ có đường cao $AE$, theo hệ thức lượng ta có:
    $SA^2 = SE \cdot SD$.
    Mặt khác, trong tam giác vuông $SAO$ tại $A$ có đường cao $AI$, theo hệ thức lượng ta có:
    $SA^2 = SI \cdot SO$.
    Từ đó suy ra $SE \cdot SD = SI \cdot SO$ (đpcm).
    \item Trong tam giác vuông $SAO$ tại $A$ có $OA = R$ và cạnh huyền $SO = 2R$, ta có:
    $\cos \widehat{AOS} = \dfrac{OA}{SO} = \dfrac{R}{2R} = \dfrac{1}{2} \implies \widehat{AOS} = 60^\circ$.
    Do tính chất tiếp tuyến đối xứng nên $\widehat{AOB} = 2\widehat{AOS} = 2 \cdot 60^\circ = 120^\circ$.
    Hình phẳng giới hạn bởi $OA, OB$ và cung nhỏ $AB$ chính là hình quạt tròn $OAB$ có góc ở tâm $120^\circ$.
    Diện tích hình quạt tròn này là:
    $S_{\text{quạt}} = \dfrac{\pi R^2 \cdot 120^\circ}{360^\circ} = \dfrac{\pi R^2}{3} = \dfrac{\pi \cdot 8^2}{3} = \dfrac{64\pi}{3} \approx 67\text{ cm}^2$.
\end{enumerate}
}
\end{bt}
